% =====================================================================
%  UML-Use-Case-Diagramm BalanceOS (UML 2.5.1, Kap. 18; OMG 2017, S. 639/641)
%  Akteur „Nutzer:in“ außerhalb der Systemgrenze „BalanceOS“,
%  Use Cases UC01–UC19 (Anforderungsdokument, Kap. 3.2)
%  Spalte A: Basis-Use-Cases (mit Assoziation zum Akteur), gruppiert nach
%            Stammdaten · Erfassung · Auswertung · Kompass · Daten · Should
%  Spalte B/C: inkludierte bzw. erweiternde Use Cases
%  Export: docs/diagramme/export/usecase.pdf
% =====================================================================
\documentclass[tikz,border=4pt]{standalone}
\usepackage{iu-diagramm}

\tikzset{
  uc/.style={usecase, text width=, minimum width=3.9cm, minimum height=1.1cm, inner ysep=0pt, inner xsep=-0.22cm},
  ucshould/.style={uc, fill=white, draw=iuGrau},
  beschr/.style={stereotyp, fill=none, inner sep=1.5pt},
}

\begin{document}
\begin{tikzpicture}

  % ---------------- Spalte A: Basis-Use-Cases ----------------------
  \def\xa{3.75}
  % Abstand innerhalb einer Gruppe 1,3 (die Ellipsen sind gut 1,1 hoch),
  % zwischen den Gruppen 1,55
  % Stammdaten
  \node[uc] (uc01) at (\xa,  0.00) {UC01\\Lebensbereiche verwalten};
  \node[uc] (uc02) at (\xa, -1.30) {UC02\\Wunschprofil festlegen};
  % Erfassung
  \node[uc] (uc03) at (\xa, -2.85) {UC03\\Aktivität erfassen};
  \node[uc] (uc04) at (\xa, -4.15) {UC04\\Aktivität bearbeiten};
  \node[uc] (uc07) at (\xa, -5.45) {UC07\\Wochenansicht anzeigen};
  % Auswertung
  \node[uc] (uc08) at (\xa, -7.00) {UC08\\Dashboard anzeigen};
  % Kompass
  \node[uc] (uc12) at (\xa, -8.55) {UC12\\Regeln konfigurieren};
  \node[uc] (uc13) at (\xa, -9.85) {UC13\\Balance-Kompass ansehen};
  % Daten
  \node[uc] (uc14) at (\xa, -11.40) {UC14\\Daten exportieren};
  \node[uc] (uc15) at (\xa, -12.70) {UC15\\Alle Daten löschen};
  \node[uc] (uc16) at (\xa, -14.00) {UC16\\Demodaten verwalten};
  % Should
  \node[ucshould] (uc17) at (\xa, -15.55) {UC17 (Should)\\WHO-5 ausfüllen};
  \node[ucshould] (uc18) at (\xa, -16.85) {UC18 (Should)\\Statistik ansehen};
  \node[ucshould] (uc19) at (\xa, -18.15) {UC19 (Should)\\Wochenreflexion erfassen};

  % ---------------- Spalte B/C: inkludierte und erweiternde UCs -----
  \def\xb{9.4}
  \node[uc] (uc05) at (\xb, -2.85) {UC05\\Wiederholung festlegen};
  \node[uc] (uc06) at (\xb, -4.15) {UC06\\Termin überspringen};
  \node[uc] (uc10) at (\xb, -7.00) {UC10\\Hinweise anzeigen};
  \node[uc] (uc11) at (\xb, -8.70) {UC11\\Begründung anzeigen};
  \node[uc] (uc09) at (\xb, -12.40) {UC09\\Wochenbewertung berechnen};

  % Kommentar (UML-Notiz) zu UC09: optionaler Gesamt-Score (US19, Should),
  % unter UC09, damit die Verbindungslinie keinen Text kreuzt
  \node[notiz, draw=none, inner xsep=4pt, anchor=north west] (n09) at (\xb-0.3, -13.55)
    {Enthält den optionalen\\Gesamt-Score mit eigener\\Gewichtung (US19, Should)};
  \draw (n09.north west) -- ([xshift=-6pt]n09.north east) -- ([yshift=-6pt]n09.north east)
        -- (n09.south east) -- (n09.south west) -- cycle;
  \draw ([xshift=-6pt]n09.north east) |- ([yshift=-6pt]n09.north east);
  \draw[dashed] (uc09.south) -- (uc09.south |- n09.north);

  % ---------------- Systemgrenze ------------------------------------
  \begin{scope}[on background layer]
    \coordinate (oben) at (\xa, 0.95);
    \node[systemgrenze, fit=(oben)(uc01)(uc19)(uc11)(n09), inner xsep=14pt, inner ysep=8pt] (sys) {};
  \end{scope}
  \node[anchor=north, font=\fontsize{10pt}{12pt}\selectfont\bfseries]
    at ([yshift=-4pt]sys.north) {BalanceOS};

  % ---------------- Akteur und Assoziationen -------------------------
  \akteur{nutzer}{-0.55,-9.1}{Nutzer:in}
  \coordinate (a) at ([xshift=0.38cm]nutzer);
  \foreach \u in {uc01,uc02,uc03,uc04,uc07,uc08,uc12,uc13,uc14,uc15,uc16,uc17,uc18,uc19}{
    \draw[assoziation] (a) -- (\u.west);
  }

  % ---------------- «extend»: erweiternder UC -> Basis-UC -----------
  \draw[extend] (uc05.west) -- node[beschr, above] {«extend»} (uc03.east);
  \draw[extend] (uc06.west) -- node[beschr, above] {«extend»} (uc04.east);
  \draw[extend] (uc11.north) -- node[beschr, right] {«extend»} (uc10.south);

  % ---------------- «include»: Basis-UC -> inkludierter UC ----------
  \draw[include] (uc08.east) -- node[beschr, above] {«include»} (uc10.west);
  \draw[include] (uc08.-22) -- node[beschr, sloped, above] {«include»} (uc09.north west);
  \draw[include] (uc18.east) -- node[beschr, sloped, above] {«include»} (uc09.south west);

  % ---------------- Legende (außerhalb der Systemgrenze) ------------
  \node[anchor=north west, align=left, font=\fontsize{8pt}{10pt}\selectfont, inner sep=0pt]
    (leg) at ([xshift=0pt, yshift=-0.35cm]sys.south west) {\textbf{Legende:}};
  \node[uc, minimum width=1.5cm, minimum height=0.55cm, anchor=west] (lm)
    at ([xshift=0.25cm]leg.east) {};
  \node[anchor=west, font=\fontsize{8pt}{10pt}\selectfont] (lmt) at (lm.east) {Must-Use-Case};
  \node[ucshould, minimum width=1.5cm, minimum height=0.55cm, anchor=west] (ls)
    at ([xshift=0.35cm]lmt.east) {};
  \node[anchor=west, font=\fontsize{8pt}{10pt}\selectfont] (lst) at (ls.east) {Should-Use-Case};
  \node[anchor=west, align=left, font=\fontsize{8pt}{10pt}\selectfont] (lit)
    at ([xshift=0.35cm]lst.east)
    {«include»: Basis-UC $\rightarrow$ inkludierter UC\\«extend»: erweiternder UC $\rightarrow$ Basis-UC};

\end{tikzpicture}
\end{document}
